\section{Data organisation and pre-processing}
\label{sec:preprocessing}

\paragraph{Organisation.} The data follows the market-basket model of RU \S6.1. The dataset is a sequence of baskets $\baskets = (T_1, \dots, T_n)$, with $n = 1000$, where each basket $T_i$ is the set of actors billed in movie $i$; identical casts of different movies remain separate baskets. In the implementation a basket is a tuple of distinct actor names sorted alphabetically, and the dataset is a list of baskets held in memory. The algorithms access the baskets only by iterating over them, once per pass, never by position, so the list could be replaced by any source that re-reads the baskets from a file at each pass, the setting assumed by RU. A one-shot stream would not suffice, because every pass after the first needs to read the baskets again.

An itemset $I$ is a set of actors, represented as a sorted tuple. Its \emph{support} is the number of movies whose basket contains it,
\[
  \supp(I) = \bigl|\{\, i \in \{1, \dots, n\} : I \subseteq T_i \,\}\bigr|,
\]
and $I$ is \emph{frequent} for a threshold $s$ if $\supp(I) \ge s$. The result of every method is a dictionary mapping each frequent itemset to its support. Sorting gives each itemset a single representation: the subsets of a sorted basket generated by \texttt{itertools.combinations} are themselves sorted, so they can be looked up directly in the dictionaries of counts and candidates, and the candidate generation of A-Priori relies on the same order.

\paragraph{Pre-processing.} Each row becomes a basket in four steps: the CSV file is read keeping empty fields as empty strings, so that no actor name is mistaken for a missing value; each of \texttt{Star1}--\texttt{Star4} is stripped of surrounding whitespace and empty names are dropped; repeated names within a movie are merged, since a basket is a set; and the remaining names are sorted.

A quality check before building the baskets found no missing, blank or whitespace-padded names. Four movies (\emph{Singin' in the Rain}, \emph{The General}, \emph{What We Do in the Shadows} and \emph{Il postino}) list the same actor in two star slots, most likely an artefact of how the source dataset was compiled; after merging, their baskets contain three actors.

\paragraph{Resulting baskets.} The result is 1000 baskets, 996 with four actors and 4 with three, over 2{,}709 distinct actors. Only 996 baskets are distinct: four pairs of sequels (\emph{The Lord of the Rings}, \emph{Star Wars} Episodes~V and~VI, \emph{Harry Potter} and \emph{Kill Bill}) have identical casts. They are kept as separate baskets, because each movie is a separate observation and support counts movies.

Actor appearances have a long tail (Table~\ref{tab:appearances}): 2{,}068 of the 2{,}709 actors appear in a single movie, and the most frequent, Robert De Niro, appears in 17. This distribution determines how many actors survive a threshold, and therefore how much work the algorithms do: lowering $s$ from 5 to 2 raises the number of frequent actors from 79 to 641.

\begin{table}[ht]
  \centering
  \caption{Number of actors by number of movies they appear in.}
  \label{tab:appearances}
  \begin{tabular}{l*{15}{r}}
    \toprule
    Movies & 1 & 2 & 3 & 4 & 5 & 6 & 7 & 8 & 9 & 10 & 11 & 12 & 13 & 14 & 17 \\
    Actors & 2068 & 370 & 126 & 66 & 20 & 29 & 13 & 2 & 6 & 1 & 3 & 2 & 1 & 1 & 1 \\
    \bottomrule
  \end{tabular}
\end{table}
